% =====================================================================
% 7. My activities in the team (PRDL, Maria Martinez)
% Table generated from docs/PLAN_DL.md by common/make_contributions.py.
% =====================================================================
\section{My activities in the team}\label{sec:activities}

The project was carried out by a team of two, and the course asks each member for an individual report
that describes her own activities. Table~\ref{tab:contrib_mm} is generated by a script from the task log
of the repository (\texttt{docs/PLAN\_DL.md}), which records who did and who reviewed each task; no entry
is typed by hand. ``Did (shared)'' marks tasks that Oriana and I did together. All code was drafted with an
AI assistant and then checked, executed and reviewed by us (Appendix~\ref{app:ai}); every run was
executed on Oriana's laptop, the runs of my blocks by me.

\begin{table}[htbp]
\centering
\caption{My role in each task of the project (generated from \texttt{docs/PLAN\_DL.md}).}
\label{tab:contrib_mm}
% Generated by common/make_contributions.py from docs/PLAN_DL.md (MM). Do not edit.
{\small
\begin{tabularx}{\textwidth}{l>{\raggedright\arraybackslash}X>{\raggedright\arraybackslash}p{3.6cm}l}
\toprule
ID & Task & Role of Mar\'ia Mart\'inez & Status \\
\midrule
T0 & Dataset shortlist (8 candidates) and choice & did (shared) & done \\
T1 & \texttt{DECISIONS.md} and \texttt{PLAN\_DL.md} first version & did (shared), reviewed & done \\
T2 & Decide the transfer-learning approach: own pre-trained model (not confirmed with the professor; raise it in the class discussion) & did (shared) & done \\
T3 & Verify licence and recommended folds on PhysioNet & did (shared) & done \\
T4 & Local project structure, \texttt{.gitignore}, environment file; GitHub repository later & reviewed & done \\
T5 & Download script, label construction, fold split, \texttt{Dataset}/\texttt{DataLoader} & did & done \\
T6 & EDA: class balance, co-occurrence, demographics, signal quality, example traces & reviewed & done \\
T7 & Evaluation module: macro AUROC, bootstrap CI, JSON logging, \texttt{make\_numbers.py} & did & done \\
T8 & Shared training loop (zero\_grad, forward, loss, backward, step; train/eval modes) & did (shared) & done \\
T9 & Dense baseline + ablations & reviewed & done \\
T10 & CNN 1D + ablations & reviewed & done \\
T11 & LSTM/GRU (+bidirectional) and RNN + attention model + ablations & did & done \\
T12 & Self-attention step by step (single layer, attention maps) and Transformer + ablations & did & done \\
T13 & Transfer learning: ECG-GPT (T23) reused on the 5 superclasses, scratch vs frozen vs fine-tuning, 10\% and 100\% labels & did (shared) & done \\
T14 & Cross-cutting ablation: 12 leads vs single lead & did (shared) & done \\
T15 & Final comparison table and figures; comparison with published results & did (shared) & done \\
T16 & Literature: dataset and CNN-related papers & reviewed & done \\
T17 & Literature: RNN/Transformer-related papers & did & done \\
T18 & Reviewed first assignment (learning strategies) for the PRDL reports & did her own part, cross review & done \\
T19 & PRDL-DL report (theory + critical analysis), opening with a summary of at most 3 pages & did her own part & done \\
T20 & MLLB-DL report (Section 0 AI environment, code, training, evaluation), opening with a summary of at most 3 pages & did her own part & done \\
T21 & \texttt{AI\_USAGE\_LOG.md} and AI Error Bounty sheets & did her own part & done \\
T22 & Adversarial review before delivery; preparation of the class discussion and oral defence & planned: owner & todo \\
T23 & Generative autoregressive model (ECG-GPT): tokenisation of lead II, causal Transformer and LSTM, sampling with temperature, ablations (D13) & reviewed & done \\
T24 & Zero-shot learning: theory and a short experiment on the class CLIP notebook & did (shared) & done \\
T25 & Initial draft of the DL reports submitted (due 1 December 2026) & planned: owner & todo \\
\bottomrule
\end{tabularx}}

\end{table}

\paragraph{My blocks.} I was responsible for the data pipeline (labels, records without a superclass,
patient-wise split, normalisation and data loaders), the evaluation protocol (metrics, F1 thresholds,
patient bootstrap, paired comparisons, run logs and the generation of the numbers of the reports), the
recurrent networks (LSTM, GRU, bidirectional networks and RNN + attention) and the self-attention and
Transformer block, including the attention module written by hand that the generative block reuses with a
causal mask. For each of them I took part in the decisions before any code, proposed the ablations, read
the code line by line with the explanation of each technique, ran the experiments, checked the results
and wrote the block note from which the reports are built. I also did the literature of the recurrent and
attention models (task T17), checked against the PDFs.

\paragraph{Shared work and review.} With Oriana I chose the dataset and the plan, wrote the training
loop, and ran the learning-strategy experiments, the zero-shot experiment, the 12-lead against lead II
ablation and the final comparison. I reviewed her project structure, exploratory analysis, dense network,
CNN, generative models and literature, and she reviewed my blocks.

\paragraph{This report.} The analysis of my blocks, the discussion and the conclusions are my own; the
summaries of Oriana's blocks are based on her block notes and attributed to her. Figures, tables and
numbers are shared by our four reports and generated from the run logs.
